\documentclass[12pt]{article}
\usepackage[brazil]{babel}
\usepackage[utf8]{inputenc}
\usepackage[T1]{fontenc}
\usepackage{setspace}
\usepackage{geometry}
\usepackage{xcolor}
\usepackage{tcolorbox}
\geometry{margin=2.5cm}

% Cores
\definecolor{azul}{RGB}{0,102,204}
\definecolor{verde}{RGB}{0,153,76}
\definecolor{roxo}{RGB}{102,0,153}
\definecolor{laranja}{RGB}{255,128,0}

\title{
\textcolor{roxo}{Roteiro de Apresentação Oral}\\
\textbf{\textcolor{azul}{Desenvolvimento de uma Mão Robótica Sustentável com Controle via Visão Computacional}}
}

\author{}
\date{}

\begin{document}

\maketitle

\onehalfspacing

\begin{tcolorbox}[colback=blue!5,colframe=azul,title=\textbf{Abertura}]
Bom dia / boa tarde.

Meu nome é \underline{\hspace{5cm}} e eu vou apresentar um projeto de \textbf{mão robótica sustentável controlada por visão computacional}.

Esse projeto mostra como a tecnologia pode ser acessível, criativa e sustentável, mesmo usando materiais simples.
\end{tcolorbox}

\begin{tcolorbox}[colback=green!5,colframe=verde,title=\textbf{Introdução}]
Com o avanço da tecnologia, a robótica passou a fazer parte do nosso dia a dia e também da educação.

Neste projeto, foi desenvolvida uma \textbf{mão robótica} construída com materiais recicláveis, como papelão, plástico reutilizado, palitos e barbante.

A mão é controlada por \textbf{visão computacional}, ou seja, ela copia os movimentos da mão humana usando uma câmera, mostrando que a robótica pode ser aprendida mesmo em contextos com poucos recursos.
\end{tcolorbox}

\begin{tcolorbox}[colback=orange!5,colframe=laranja,title=\textbf{Objetivo do Projeto}]
O objetivo principal do projeto é desenvolver uma \textbf{mão robótica funcional}, feita com materiais recicláveis e controlada por visão computacional.

Além disso, o projeto busca promover a aprendizagem prática, a sustentabilidade e o uso da tecnologia para fins educacionais e assistivos.
\end{tcolorbox}

\begin{tcolorbox}[colback=purple!5,colframe=roxo,title=\textbf{Como o Projeto Funciona}]
O funcionamento da mão robótica acontece em três etapas principais:

\textbf{1º passo -- Visão computacional:} Uma câmera captura os movimentos da mão humana, e um programa em \textbf{Python} reconhece gestos como abrir e fechar os dedos.

\textbf{2º passo -- Integração eletrônica:} Os comandos identificados pelo programa são enviados para o \textbf{Arduino}, que controla os servomotores.

\textbf{3º passo -- Estrutura mecânica sustentável:} A mão robótica é construída com materiais recicláveis, que se movimentam por meio de fios e motores.
\end{tcolorbox}

\begin{tcolorbox}[colback=orange!5,colframe=laranja,title=\textbf{Tecnologias Utilizadas}]
\textbf{Visão Computacional:} Responsável por identificar os gestos da mão humana através da câmera.

\textbf{Arduino:} Controla os servomotores e executa os comandos recebidos do sistema.

\textbf{Python:} Linguagem usada para processar as imagens e enviar os comandos para o Arduino.

\textbf{Materiais Recicláveis:} Papelão, plástico reutilizado, palitos e barbante formam a estrutura da mão robótica.
\end{tcolorbox}

\begin{tcolorbox}[colback=blue!5,colframe=azul,title=\textbf{Resultados}]
Durante os testes, a mão robótica conseguiu reproduzir movimentos básicos, como abrir e fechar os dedos.

O desempenho foi satisfatório, mas a precisão pode variar conforme a iluminação do ambiente e a qualidade da câmera.
\end{tcolorbox}

\begin{tcolorbox}[colback=green!5,colframe=verde,title=\textbf{Conclusão}]
Com esse projeto, foi possível perceber que a robótica pode ser \textbf{acessível, educativa e sustentável}.

A integração entre visão computacional, Arduino e materiais recicláveis mostrou grande potencial para o ensino de tecnologia e para futuras aplicações assistivas.

Esse trabalho reforça que é possível aprender robótica de forma prática, criativa e consciente.

\textbf{Muito obrigado pela atenção!}
\end{tcolorbox}

\end{document}

